% الجزء: sets-functions-relations
% الفصل: sets
% القسم: important-sets

\documentclass[../../../include/open-logic-section]{subfiles}

\begin{document}

\olfileid[ar]{sfr}{set}{imp}
\olsection{بعض المجموعات المهمة}

\begin{ex}
سنتعامل في الغالب مع مجموعات تكون !!{element}s كل منها كائنات
رياضية. ومنها أربع مجموعات مهمة بما يكفي لأن تكون لها
أسماء مخصوصة:
\begin{multline*}
    \Nat = \{0, 1, 2, 3, \ldots\} \\
    \shoveright{\text{مجموعة الأعداد الطبيعية}}\\
    \shoveleft{\Int = \{\ldots, -2, -1, 0, 1, 2, \ldots\}} \\
    \shoveright{\text{مجموعة الأعداد الصحيحة}}\\
    \shoveleft{\Rat = \Setabs{\nicefrac{m}{n}}{m, n \in \Int\text{ و }n \neq 0}}\\
    \shoveright{\text{مجموعة الأعداد النسبية}}\\
    \shoveleft{\Real = (-\infty, \infty)}\\
    \text{مجموعة الأعداد الحقيقية (المتصل العددي)}
\end{multline*}
وهذه كلها مجموعات \emph{لا نهائية}، أي إن لكل منها عددًا
لا نهائيًا من !!{element}s.

وعندما ننتقل من مجموعة إلى التي تليها في هذه القائمة، نضيف
\emph{مزيدًا} من الأعداد إلى ما لدينا. وبالفعل، من الواضح أن
$\Nat \subseteq \Int \subseteq \Rat \subseteq \Real$؛ فكل عدد طبيعي
عدد صحيح، وكل عدد صحيح عدد نسبي، وكل عدد نسبي عدد حقيقي.
ومن الواضح بالقدر نفسه أن $\Nat \subsetneq \Int \subsetneq
\Rat$، لأن $-1$ عدد صحيح لكنه ليس عددًا طبيعيًا، ولأن
$\nicefrac{1}{2}$ عدد نسبي لكنه ليس عددًا صحيحًا. وأقل وضوحًا
أن $\Rat \subsetneq \Real$، أي إن هناك أعدادًا حقيقية غير
نسبية\oliflabeldef{sfr:arith:real:realline}{، لكننا سنعود إلى ذلك في
\olref[arith][real]{realline}}{}.

وسنستخدم أحيانًا أيضًا مجموعة الأعداد الصحيحة الموجبة $\PosInt = \{1,
2, 3, \dots\}$، والمجموعة التي لا تضم إلا أول عددين طبيعيين
$\Bin = \{0, 1\}$.
\end{ex}

\begin{tagblock}{compsci}
\begin{ex}[السلاسل]
مثال آخر جدير بالاهتمام هو مجموعة \emph{السلاسل المنتهية} $A^{*}$
على أبجدية $A$: فكل متتالية منتهية من عناصر~$A$ سلسلة على $A$.
وندرج \emph{السلسلة الفارغة $\Lambda$} ضمن السلاسل على~$A$ لكل
أبجدية~$A$. فمثلًا،
\begin{multline*}
\Bin^*
=\{\Lambda,0,1,00,01,10,11,\\
000,001,010,011,100,101,110,111,0000,\ldots\}.
\end{multline*}
إذا كانت $x=x_{1}\ldots x_{n}\in A^{*}$ سلسلة مؤلفة من $n$
«حروف» من $A$، فإننا نقول إن \emph{طول} السلسلة هو~$n$،
ونكتب $\len{x}=n$.
\end{ex}
\end{tagblock}

\begin{ex}[المتتاليات اللانهائية]
ولكل مجموعة $A$ يمكننا أيضًا النظر في المجموعة~$A^\omega$ المؤلفة
من متتاليات لا نهائية من !!{element}s المجموعة~$A$. وتتكون المتتالية
اللانهائية $a_1a_2a_3a_4\dots$ من قائمة لا نهائية أحادية الاتجاه من
كائنات، كل واحد منها !!a{element} من~$A$.
\end{ex}

\end{document}
